\emph{Orientation.} By Fact \claimref{odd-extensions.imported-conformal-completion} the conformal ambient of a representative of signature $(p,q)$ has signature $(p+1, q+1)$ and dimension $n + 2$. If $n$ is odd the ambient dimension is odd and Proposition \claimref{dark-matter.two-families-are-two-chiralities} gives a single family, so no Witt pair lies in opposite families and the tower is not oriented. If $n$ is even the ambient has dimension $2m$ with $m = (n+2)/2$, and by part (iii) of that proposition a Witt pair lies in opposite families exactly when $m$ is odd, that is $n + 2 \equiv 2 \pmod 4$, that is $n \equiv 0 \pmod 4$. So a tower is oriented if and only if $n \equiv 0 \pmod 4$, and in particular $n$ is even, and so is $r$, since $p - q \equiv p + q \pmod 2$.

\emph{(i).} Let the tower be oriented and causal. Causality, $\min(p,q) = 1$, needs $n \ge 2$, and with $n \equiv 0 \pmod 4$ this gives $n \ge 4$. Then the signatures of dimension $n$ with $\min(p,q) = 1$ are $(1, n-1)$ and $(n-1, 1)$, which are distinct because $n - 1 \ge 3 > 1$; there are exactly two. Their residues are $1 - (n-1) = 2 - n$ and $(n-1) - 1 = n - 2$. From $n \equiv 0 \pmod 4$, $n - 2 \equiv 2 \pmod 4$, so $n - 2 \equiv 2$ or $6 \pmod 8$, and $2 - n = -(n-2)$ is then $6$ or $2$ respectively: the two residues are $\{2, 6\}$ as a set, and they are exchanged by $r \mapsto -r$. Every oriented causal tower is one of these two, so has residue $2$ or $6$.

\emph{(ii).} Modulo $8$ the map $r \mapsto -r$ fixes exactly $0$ and $4$ ($2r \equiv 0 \pmod 8$ iff $r \equiv 0 \pmod 4$), and by (i) neither is the residue of an oriented causal tower, in any dimension. For any signature with $r \equiv 0$ or $4 \pmod 8$, $p - q \equiv 0 \pmod 4$, so Lemma \claimref{transport.sign} gives $\omega^{2} = +1$. An element squaring to $+1$ is not an imaginary unit, so no tower of these residues has a pseudoscalar imaginary unit. Here $n \equiv r \pmod 2$ is even, so $\omega$ has even grade; a frame vector $e_{i}$ anticommutes with the $n - 1$ other factors of $\omega = e_{1} \cdots e_{n}$ and commutes with itself, so $e_{i}\omega = (-1)^{n-1}\omega e_{i}$, and for $n$ even every vector anticommutes with $\omega$; hence every product of two vectors commutes with $\omega$, and $\omega$ is central in the even subalgebra. With $\omega^{2} = 1$, $\bigl(\tfrac12(1 \pm \omega)\bigr)^{2} = \tfrac14(2 \pm 2\omega) = \tfrac12(1 \pm \omega)$, the two idempotents are central in the even subalgebra, sum to $1$ and multiply to $\tfrac14(1 - \omega^{2}) = 0$: chirality is a splitting into two real ideals. (For $n = 0$, the signature $(0,0)$, $\omega = 1$ and the splitting is trivial, $1 = 1 + 0$; for $n \ge 2$ both idempotents are nonzero, since $\omega \notin \R$.)

\emph{(iii).} An oriented causal tower has $r = \pm(n - 2) \equiv 2 \pmod 4$, so Lemma \claimref{transport.sign} gives $\omega_{n}^{2} = -1$. The conformal ambient $(p+1, q+1)$ has the same $p - q$, so the same lemma gives $\omega_{n+2}^{2} = -1$. Since $n$ is even, $\omega_{n}$ anticommutes with each of the tower's vectors (the argument of (ii)) and is central in the even subalgebra. A vector of the ascended plane is orthogonal to all $n$ factors of $\omega_{n}$ and anticommutes with each, an even number of sign changes, so it commutes with $\omega_{n}$. These are the properties Proposition \claimref{field-theory.where-each-lives} records for $I_{4}$, which at $n = 4$, $(p,q) = (1,3)$, is $\omega_{4}$ itself: $I_{4}$ spans the centre of $\Cl^{+}(1,3)$ with $1$, and $[I_{4}, E] = 0$ for the ascended plane's bivector $E$. Finally $p - q \equiv 2 \pmod 4$ and $n$ is even, so by Theorem \claimref{rule2.where-obstruction-sits-commuting-square}(a) the spacelike extension $\Cl(p, q+1)$ carries the type-C obstruction. Its volume element $\omega_{n+1}$ has odd grade $n + 1$, so it commutes with every vector and is central, and $p - (q+1) = r - 1 \equiv 1 \pmod 4$ gives $\omega_{n+1}^{2} = +1$ by Lemma \claimref{transport.sign}; so $\tfrac12(1 \pm \omega_{n+1})$ are central idempotents, nonzero because $\omega_{n+1} \notin \R$, summing to $1$, and they are the two ideals of the obstructed extension.

\emph{The consequence.} The argument of (i) used only $n \equiv 0 \pmod 4$ and $\min(p,q) = 1$ and made no choice of dimension, so it holds for every representative of every residue. An odd residue has $n$ odd and is never oriented; residues $0$ and $4$ are excluded by (ii); so in every window the oriented causal towers have residues $2$ and $6$ only, which is the conclusion Corollary \claimref{dark-matter.which-towers-can-carry-arrow} reached at the four-dimensional representatives of Premise P3, now without that premise.

The check computes, for every signature of dimension at most $10$, the family bit from $m = (n+2)/2$, causality, the residue and $\omega^{2}$ from the blade product, and confirms (i) and (ii) by enumeration; for the oriented causal signatures of dimensions $4$ and $8$ it verifies in exact arithmetic that $\omega_{n}$ anticommutes with every tower vector, commutes with both ascended generators and every bivector, that $\omega_{n+2}^{2} = -1$, and that in $\Cl(p, q+1)$ the volume element is central with square $+1$.
